\chapter{Implementation}
\label{chapter_implementation}

\section{Implementation Overview}

The proposed work will be implemented as a simulation-based optimization
study rather than as a physical UAV deployment. The implementation will use a
Python-based experimental workflow that can be executed locally and later
shared through a reproducible notebook environment such as Google Colab or
Jupyter Notebook. This choice keeps the experimental setup transparent while
allowing the task-allocation, routing, payload, battery, and recourse
components of the proposed PGBM formulation to be tested under controlled
post-earthquake response scenarios.

The implementation will focus on generating and solving representative
responder-UAV planning instances. Each instance will contain an earthquake-
affected urban workspace, accepted survivor-assistance tasks, available
responder UAVs, relief-item requirements, feasible three-dimensional flight
links, and possible newly arriving tasks for recourse evaluation. The purpose
of this chapter is therefore to describe the planned data and environment
setup that will be used to evaluate the mathematical formulation presented in
Chapter~\ref{chapter_proposedMethod}.

\section{Data Sources and Scenario Construction}

Public disaster-response datasets will be used to support the scenario design,
but they will not be treated as complete optimization datasets. DRespNeT is a
UAV-based post-earthquake dataset containing high-resolution aerial imagery
with detailed annotations of damaged structures, debris, access points, rescue
personnel, vehicles, and visible civilians \cite{sirma2025drespnet}. It
provides a realistic reference for understanding post-earthquake urban
environments and UAV-assisted search-and-rescue scenarios.

LADI, the Low Altitude Disaster Imagery dataset, contains aerial images
collected from real disaster-affected regions and annotated for disaster-
related visual features \cite{scheele2025ladi}. It will be used as a reference
for the visual characteristics and spatial complexity of emergency
environments observed from low-altitude aerial platforms. AIDER, the Aerial
Image Dataset for Emergency Response, contains aerial images representing
emergency conditions such as collapsed buildings, rubble, fires, floods,
traffic incidents, and normal scenes \cite{kyrkou2020emergencynet}. It is
useful for supporting disaster-scene classification and emergency-response
scenario design.

These public datasets are used primarily to support the design and realism of
the disaster scenario. The task, UAV, payload, energy, and dynamic recourse
data required by the proposed optimization model will be generated
synthetically, because the available public datasets do not jointly provide
survivor-assistance demand, responder-UAV battery state, payload state,
obstacle-aware three-dimensional routing links, and recourse-event records.

\section{Synthetic Survivor-Assistance Task Generation}

The synthetic task generator will create survivor-assistance tasks that match
the inputs required by the proposed formulation. Each task will include a task
identifier, detected survivor location, UAV-accessible drop-off waypoint,
detection time, operational severity score, required relief-item quantities,
and service duration. The detected survivor location and the drop-off waypoint
will be stored separately, since a survivor may be located inside rubble or a
damaged structure while the UAV can only deliver from a nearby feasible
position in flyable space.

Task arrival times will also be generated so that both initial planning and
dynamic recourse can be tested. Tasks known at the beginning of a planning
cycle will form the initial task set, while later arrivals will be used to
evaluate whether the task-replacement rule can improve the remaining mission
value without restarting the entire plan.

\section{Three-Dimensional Operating Environment}

The main experimental environment will be a custom synthetic three-dimensional
urban post-earthquake workspace. The workspace will be represented as a bounded
3D region containing UAV base locations, building obstacles, rubble-affected
areas, and no-fly volumes. Feasible UAV waypoints will be sampled only from
flyable space, and directed flight links will be created only when the
straight-line segment between two waypoints does not intersect a blocked
region.

This environment is consistent with simulation-based UAV disaster-response
studies, where researchers construct 3D disaster rescue scenarios, terrain
models, threat sources, and mission points to evaluate assignment and path-
planning methods \cite{xiong2023multidrone,liu2021multiuv}. For the proposed
PGBM implementation, the environment will be kept simple enough to support
reproducible experiments, while still preserving the key operational
constraints: obstacle avoidance, feasible drop-off points, limited battery,
limited payload, and return to the same physical UAV base.

\section{Responder-UAV and Payload Configuration}

Each responder UAV will be described by its base location, cruise speed,
available battery energy, reserve energy requirement, maximum payload weight,
maximum parcel count, and energy-related coefficients. Relief supplies will be
represented using a finite set of item types, where each item type has a
specified unit weight. The generated task demand will determine how many items
of each type must be loaded for a selected task.

The same UAV and payload configuration will be used consistently across the
assignment, routing, timing, payload, and energy calculations. In particular,
each planned mission must start from the UAV base, visit the assigned
drop-off waypoints through feasible flight links, and return to the same
physical base while preserving the required reserve energy.



The implementation workflow will begin by generating a disaster scenario and
constructing the corresponding 3D flight-link graph. The planner will then
create the current set of eligible survivor-assistance tasks and the available
responder-UAV states. The optimization model will be solved on this generated
instance to obtain the selected tasks, assigned UAVs, selected flight links,
completion times, payload usage, and energy-feasible missions.

During simulated mission execution, the state of each UAV and each task will
be updated. If a new survivor-assistance task arrives while responder UAVs are
already active, the recourse module will evaluate whether replacing one
uncompleted task in an active mission improves the remaining service value.
Accepted replacements will update the remaining mission plan, while rejected
new tasks or displaced old tasks will remain in the waiting queue for a later
planning cycle.

\section{Planned Evaluation Metrics}

The planned experiments will evaluate both the initial planning decision and
the dynamic recourse behavior. The main output measures will include the total
time-dependent service value, number of served tasks, average completion
delay, total UAV travel distance, battery-energy usage, payload utilization,
number of deferred tasks, number of accepted task replacements, and number of
rejected recourse events. Safe-return feasibility will also be checked for
each UAV mission.

\begin{table}[htbp]
\centering
\caption{Planned implementation data components.}
\label{tab:implementation_components}
\begin{tabular}{p{0.28\textwidth}p{0.58\textwidth}}
\hline
\textbf{Component} & \textbf{Planned Data}\\
\hline
Scenario references & DRespNeT, LADI, and AIDER for disaster-scene context\\
Task instances & Survivor location, drop-off waypoint, severity, demand, and service time\\
UAV configuration & Base, speed, payload capacity, parcel capacity, battery, and reserve energy\\
3D environment & Workspace boundary, obstacles, feasible waypoints, and flight links\\
Recourse events & Newly arriving tasks during active responder-UAV missions\\
Evaluation outputs & Service value, delay, distance, energy use, payload use, and replacement decisions\\
\hline
\end{tabular}
\end{table}
